\section{Tangent measurement on the archived spike-region history}
\label{sec:archive}
The recovered finite-window history used the archived
CBE3b data at $\tau=18,16$ with $\tau=-\log t$.
To check conditioning at this trajectory, restart the centered
second-order spatial method at the recorded $\tau=18$ data and
measure
\[
 J=-\frac12 \frac{\partial P(16)}{\partial P_\tau(18)}
      \bigg|_{P(18),Q(18),Q_\tau(18)} .
\]
The normalization gives the identity for free motion in $\tau$.
The linearization was integrated jointly with the replay.
Its nearest-neighbor structure in the variables
$(\delta P,\delta P_\tau,e^P\delta Q,e^P\delta Q_\tau)$
allowed periodic coloring of tangent seeds.

The computed maximum row defects from the identity at
$N=1024,2048,4096$ were $0.0619925$, $0.0620070$, and
$0.0620106$. At the finest resolution the diagonal ran
from $0.937989$ to $1.000000$ to the shown precision, and the
estimated maximum off-diagonal row sum was $1.42\times10^{-9}$.
Refining both the time step ceiling and the number of colors
changed the maximum defect by less than $7\times10^{-14}$.
A centered finite difference in a fixed mixed Fourier
direction matched the tangent evolution to within
$8.7\times10^{-10}$ at all three resolutions. The terminal
$P$ replay matched the archival value within $8.2\times10^{-12}$.

To track possible color aliasing, suppose the exact scaled
generator has uniform in-time bounds $D$ and $E$ on its
block-diagonal and adjacent-site parts in maximum row norm.
The block-diagonal propagator preserves support and each
adjacent-site insertion moves support by at most one site.
For an interval of duration $\Delta$, the norm of transfer
outside radius $m$ is thus at most
\[
 e^{\Delta(D+E)}
   \frac{(\Delta E)^{m+1}}{(m+1)!}.
\]
Apply the same factor $1/\Delta$ used in $J$ to its projected
tail. The sampled $D,E$ at $N=4096$ give a diagnostic tail
scale of $3.05\times10^{-17}$ for radius three. These
generator maxima were sampled in time. They and the tangent
results are numerical evidence at the stated discretizations.
An independently verified bound of $0.0621$ on the full defect
would imply $\|J^{-1}\|_\infty\leq 1.067$ there by the
Neumann series. The observed small defect provides a local
explanation for the successful nonlinear restart and is
consistent with the exact failures at other observation windows.
